% Source: Weinberg GR, physical PDF pp. 123--125 (printed pp. 98--100).
% Coverage: Section 4.4, Tensor Densities; equations (4.4.1)--(4.4.11).
% Figures/tables/footnotes: one linked reference marker; no figures or tables.
% Status: source-reviewed and compile-clean.
% Uncertainties: none.

\subsection{Tensor Densities}\label{sec:4.4}

Despite the ubiquity of tensors, there is nothing sacred about the tensor
transformation law. One very important example of a nontensor is the
determinant of the metric tensor
\begin{equation}
  g\equiv\det(g_{\mu\nu})<0.
  \tag{4.4.1}\label{eq:4.4.1}
\end{equation}
The transformation rule for the metric tensor can be regarded as a matrix
equation,
\[
  g'_{\mu\nu}
  =
  \frac{\partial x^\rho}{\partial x'^\mu}
  g_{\rho\sigma}
  \frac{\partial x^\sigma}{\partial x'^\nu},
\]
and taking the determinant, we find that
\begin{equation}
  g'
  =
  \left\lvert\frac{\partial x}{\partial x'}\right\rvert^2 g,
  \tag{4.4.2}\label{eq:4.4.2}
\end{equation}
where \(\lvert\partial x/\partial x'\rvert\) is the Jacobian of the
transformation \(x'\to x\), that is, it is the determinant of the matrix
\(\partial x^\mu/\partial x'^\nu\). A quantity such as \(g\), which
transforms like a scalar except for extra factors of the Jacobian, is called
a \textit{scalar density}, and similarly a quantity that transforms as a
tensor except for extra factors of the Jacobian determinant is called a
\textit{tensor density}. The number of factors of the determinant
\(\lvert\partial x'/\partial x\rvert\) is called the \textit{weight} of the
density; for instance, we see from Eq.~\eqref{eq:4.4.2} that \(g\) is a
density of weight \(-2\), because
\begin{equation}
  \left\lvert\frac{\partial x}{\partial x'}\right\rvert
  =
  \left\lvert\frac{\partial x'}{\partial x}\right\rvert^{-1},
  \tag{4.4.3}\label{eq:4.4.3}
\end{equation}
as can be seen by taking the determinant of the equation
\[
  \frac{\partial x^\mu}{\partial x'^\lambda}
  \frac{\partial x'^\lambda}{\partial x^\nu}
  =\delta^\mu{}_\nu.
\]

Any tensor density of weight \(W\) can be expressed as an ordinary tensor
times a factor \((-g)^{-W/2}\). For instance, a tensor density
\(\mathcal T^\mu{}_\nu\) of weight \(W\) has the transformation rule
\begin{equation}
  \mathcal T'^\mu{}_\nu
  =
  \left\lvert\frac{\partial x'}{\partial x}\right\rvert^W
  \frac{\partial x'^\mu}{\partial x^\lambda}
  \frac{\partial x^\kappa}{\partial x'^\nu}
  \mathcal T^\lambda{}_\kappa,
  \tag{4.4.4}\label{eq:4.4.4}
\end{equation}
and using Eq.~\eqref{eq:4.4.2}, we see that
\begin{equation}
  (-g')^{W/2}\mathcal T'^\mu{}_\nu
  =
  \frac{\partial x'^\mu}{\partial x^\lambda}
  \frac{\partial x^\kappa}{\partial x'^\nu}
  (-g)^{W/2}\mathcal T^\lambda{}_\kappa.
  \tag{4.4.5}\label{eq:4.4.5}
\end{equation}

The importance of tensor densities arises from the fundamental theorem of
integral calculus,~\hyperref[ch4-ref-4]{[4]} that under a general coordinate
transformation \(x\to x'\), the volume element \(\dd^4x\) becomes
\begin{equation}
  \dd^4x'
  =
  \left\lvert\frac{\partial x'}{\partial x}\right\rvert\dd^4x.
  \tag{4.4.6}\label{eq:4.4.6}
\end{equation}
Hence the product of \(\dd^4x\) with a tensor density of weight \(-1\)
transforms like an ordinary tensor. In particular,
\(\sqrt{-g}\,\dd^4x\) is an invariant volume element.

There is one tensor density whose components are the same in all
orientation-preserving coordinate systems: the permutation symbol
\(\tilde\epsilon^{\mu\nu\lambda\kappa}\). To define it, we order the
coordinate indices in the reference sequence \(0,1,2,3\), corresponding to
\(t,x,y,z\). Then \(\tilde\epsilon^{\mu\nu\lambda\kappa}\) is defined by
\begin{equation}
  \tilde\epsilon^{\mu\nu\lambda\kappa}
  =
  \begin{cases}
    +1,& \mu\nu\lambda\kappa\text{ an even permutation of }0123,\\
    -1,& \mu\nu\lambda\kappa\text{ an odd permutation of }0123,\\
    0,& \text{some indices equal}.
  \end{cases}
  \tag{4.4.7}\label{eq:4.4.7}
\end{equation}
To see that this is a tensor density, consider the quantity
\begin{equation}
  \frac{\partial x'^\rho}{\partial x^\mu}
  \frac{\partial x'^\sigma}{\partial x^\nu}
  \frac{\partial x'^\eta}{\partial x^\lambda}
  \frac{\partial x'^\zeta}{\partial x^\kappa}
  \tilde\epsilon^{\mu\nu\lambda\kappa}.
  \tag{4.4.8}\label{eq:4.4.8}
\end{equation}
We note that this is completely antisymmetric in the indices
\(\rho,\sigma,\eta,\zeta\), and therefore proportional to
\(\tilde\epsilon^{\rho\sigma\eta\zeta}\). To determine the proportionality
constant, let \(\rho\sigma\eta\zeta\) take the values of the reference
sequence. For an orientation-preserving transformation, define the positive
Jacobian \(J\equiv\det(\partial x'/\partial x)>0\). Then
Eq.~\eqref{eq:4.4.8} is just \(J\), and so
\begin{equation}
  \frac{\partial x'^\rho}{\partial x^\mu}
  \frac{\partial x'^\sigma}{\partial x^\nu}
  \frac{\partial x'^\eta}{\partial x^\lambda}
  \frac{\partial x'^\zeta}{\partial x^\kappa}
  \tilde\epsilon^{\mu\nu\lambda\kappa}
  =
  J\tilde\epsilon^{\rho\sigma\eta\zeta}.
  \tag{4.4.9}\label{eq:4.4.9}
\end{equation}
Thus \(\tilde\epsilon^{\mu\nu\lambda\kappa}\) is a contravariant tensor
density of weight \(-1\) for orientation-preserving maps. Under an
orientation reversal it acquires the additional sign characteristic of a
pseudodensity. The ordinary contravariant Levi-Civita tensor is
\[
  \varepsilon^{\mu\nu\lambda\kappa}
  \equiv
  \frac{1}{\sqrt{-g}}\,
  \tilde\epsilon^{\mu\nu\lambda\kappa}.
\]
We lower its indices in the usual tensorial way:
\begin{equation}
  \varepsilon_{\rho\sigma\eta\zeta}
  \equiv
  g_{\rho\mu}g_{\sigma\nu}g_{\eta\lambda}g_{\zeta\kappa}
  \varepsilon^{\mu\nu\lambda\kappa}.
  \tag{4.4.10}\label{eq:4.4.10}
\end{equation}
If the lower-index permutation symbol is separately normalized by
\(\tilde\epsilon_{0123}=+1\), the mostly-plus signature gives
\begin{equation}
  \varepsilon_{\rho\sigma\eta\zeta}
  =
  -\sqrt{-g}\,
  \tilde\epsilon_{\rho\sigma\eta\zeta}.
  \tag{4.4.11}\label{eq:4.4.11}
\end{equation}
In particular,
\(\varepsilon^{0123}=+1/\sqrt{-g}\) and
\(\varepsilon_{0123}=-\sqrt{-g}\).

The rules of tensor algebra may be easily extended to encompass tensor
densities:

\begin{enumerate}
  \item[(A)] The linear combination of two tensor densities of the
  \textit{same} weight \(W\) is a tensor density of weight \(W\).

  \item[(B)] The direct product of two tensor densities of weights \(W_1\)
  and \(W_2\) yields a tensor density of weight \(W_1+W_2\).

  \item[(C)] The contraction of indices on a tensor density of weight \(W\)
  yields a tensor density of weight \(W\). From (B) and (C) it follows that
  raising and lowering indices does not change the weight of a tensor
  density.
\end{enumerate}
